\chapter{Lecture 4 --- Networking Devices, Topologies and the Cloud Network}

\begin{keyidea}[The one idea from this lecture]
Every networking device is best described by \textbf{which layer it works at}, i.e.\ \textbf{how deep into the packet it is willing to look}. A hub looks at nothing, a switch looks at MAC addresses (L2), a router looks at IP addresses (L3), a firewall looks at IPs and ports (L3/L4). Then we decide how to \textbf{arrange} the devices (topology), and see that the cloud does the same thing in \textbf{software}.
\end{keyidea}

\section{Why do we need network devices?}

\subsection{Connecting everyone directly does not scale}
If every computer gets its own cable to every other computer (a full mesh), the number of cables for $n$ devices is:
\[ \text{cables} = \frac{n(n-1)}{2} \]

\begin{center}
\begin{tabular}{rlr}
\toprule
Devices & Calculation & Cables needed \\
\midrule
3 & $3\times 2/2$ & 3 \\
10 & $10\times 9/2$ & 45 \\
100 & $100\times 99/2$ & 4{,}950 \\
1{,}000 & $1000\times 999/2$ & 499{,}500 \\
\bottomrule
\end{tabular}
\end{center}
Devices grow \textbf{linearly}, cables grow with the \textbf{square}. Ten times the devices $\approx$ a hundred times the cabling. Impossible for an office, let alone a data centre.

\subsection{The pivot: everyone connects to ONE device}
\begin{center}
\begin{tikzpicture}
  \begin{scope}
    \foreach \a [count=\i] in {90,150,210,270,330,30} \node[host,minimum width=8mm,minimum height=6mm] (m\i) at (\a:1.6) {};
    \foreach \i in {1,...,6} \foreach \j in {1,...,6} { \ifnum\i<\j \draw[gray] (m\i) -- (m\j); \fi }
    \node[below=1mm,font=\small] at (0,-2.0) {Full mesh: $6\cdot5/2=15$ cables};
  \end{scope}
  \begin{scope}[xshift=6.5cm]
    \node[draw,fill=cnorange!25,minimum width=12mm,minimum height=7mm,font=\small] (c) at (0,0) {device};
    \foreach \a [count=\i] in {90,150,210,270,330,30} { \node[host,minimum width=8mm,minimum height=6mm] (s\i) at (\a:1.7) {}; \draw[thick] (c) -- (s\i); }
    \node[below=1mm,font=\small] at (0,-2.0) {Central device: just 6 cables};
  \end{scope}
\end{tikzpicture}
\end{center}
Now $n$ devices need only $n$ cables. \textbf{This single move is why networking hardware exists.} The rest of the lecture asks: what should that central device actually be?

\section{Networking devices}
Core components that connect computers and other devices: \textbf{hub, switch, wireless router, access point}, plus \textbf{firewall} and \textbf{load balancer}.

\subsection{The hub --- the 1980s answer (Layer 1)}
\begin{itemize}
  \item \textbf{Broadcasts} incoming data to \textbf{all} connected devices.
  \item Does \textbf{not} recognise hardware (MAC) addresses at all.
  \item Works at the \textbf{Physical layer (L1)}: takes the electrical signal on one port and repeats it out of every other port.
\end{itemize}

\begin{center}
\begin{tikzpicture}
  \node[draw,fill=gray!25,minimum width=14mm,minimum height=7mm] (h) at (0,0) {HUB};
  \node[host] (a) at (-3,0) {A};
  \node[host] (b) at (3,1) {B};
  \node[host] (c) at (3,0) {C};
  \node[host] (d) at (3,-1) {D};
  \draw[arr,cnblue] (a) -- node[above,font=\scriptsize]{to C} (h);
  \draw[arr,cnred] (h) -- (b); \draw[arr,cngreen] (h) -- (c); \draw[arr,cnred] (h) -- (d);
  \node[font=\scriptsize,cnred,align=left] at (5.6,0.5) {B and D get it too,\\then discard it};
\end{tikzpicture}
\end{center}

\begin{examtip}[Class question: main real-world problem with a hub?]
\textbf{Unnecessary traffic.} Data goes to every device even when only one needs it. Every machine must receive, inspect and throw away traffic not meant for it, and while that traffic occupies the shared medium, nobody else can transmit. More devices $\Rightarrow$ worse congestion.
\end{examtip}

\subsection{The switch --- the 1990s answer (Layer 2)}
\begin{itemize}
  \item Works at the \textbf{Data Link layer (L2)}.
  \item Connects devices within a \textbf{LAN} and forwards data to the \textbf{correct} device using \textbf{MAC addresses}.
  \item Has \textbf{no understanding of Layer 3} --- IP addresses mean nothing to it.
  \item Ports can be configured; broadcasting can be controlled because it knows which MAC sits on which port.
\end{itemize}
\textbf{``The switch is a hub with a memory.''} It learns which MAC address lives behind which port (its MAC table), so it sends each frame only where it needs to go.

\begin{center}
\begin{tikzpicture}
  \node[draw,fill=cnblue!20,minimum width=16mm,minimum height=7mm] (s) at (0,0) {SWITCH};
  \node[host] (a) at (-3,0) {A};
  \node[host] (b) at (3,1) {B};
  \node[host] (c) at (3,0) {C};
  \node[host] (d) at (3,-1) {D};
  \draw[arr,cnblue] (a) -- node[above,font=\scriptsize]{to C} (s);
  \draw[gray] (s) -- (b); \draw[arr,cngreen] (s) -- (c); \draw[gray] (s) -- (d);
  \node[draw,font=\scriptsize,below=6mm of s,align=left] {\textbf{MAC table}\\MAC-A $\to$ port 1\\MAC-B $\to$ port 2\\MAC-C $\to$ port 3\\MAC-D $\to$ port 4};
\end{tikzpicture}
\end{center}

\subsection{Hub vs switch}
\begin{center}
\begin{tabularx}{\textwidth}{>{\bfseries}l X X}
\toprule
 & Hub & Switch \\
\midrule
Era & 1980s & 1990s onward \\
Layer & L1 --- Physical & L2 --- Data Link \\
Addresses understood & None & MAC addresses \\
Forwarding & Broadcasts to every port & Sends to the correct port only \\
Main weakness & Unnecessary traffic; congestion grows with devices & Cannot see IP, so cannot reach another network \\
\bottomrule
\end{tabularx}
\end{center}

\subsection{The router --- reaching another network (Layer 3)}
A switch solves communication \emph{inside} one network. \textbf{What if we want to talk to a computer in another network?} The switch can't help (no L3 knowledge). The \textbf{router} can:
\begin{itemize}
  \item Routes data \textbf{between different networks}.
  \item Works at the \textbf{Network layer (L3)} using \textbf{IP addresses}. (The GPS from Lecture 3.)
\end{itemize}

\begin{center}
\begin{tikzpicture}
  \node[draw,dashed,rounded corners,minimum width=40mm,minimum height=26mm] (n1) at (0,0) {};
  \node[font=\footnotesize] at (0,1.0) {Network 1 (LAN)};
  \node[draw,fill=cnblue!20,font=\scriptsize] (s1) at (0,0) {Switch};
  \node[host,minimum width=8mm,font=\scriptsize] at (-1.3,-0.8) {PC}; \node[host,minimum width=8mm,font=\scriptsize] at (1.3,-0.8) {PC};
  \draw (s1) -- (-1.3,-0.5); \draw (s1) -- (1.3,-0.5);
  \node[draw,dashed,rounded corners,minimum width=40mm,minimum height=26mm] (n2) at (8,0) {};
  \node[font=\footnotesize] at (8,1.0) {Network 2 (LAN)};
  \node[draw,fill=cnblue!20,font=\scriptsize] (s2) at (8,0) {Switch};
  \node[host,minimum width=8mm,font=\scriptsize] at (6.7,-0.8) {PC}; \node[host,minimum width=8mm,font=\scriptsize] at (9.3,-0.8) {PC};
  \draw (s2) -- (6.7,-0.5); \draw (s2) -- (9.3,-0.5);
  \node[rtr,minimum size=12mm] (r) at (4,0) {Router};
  \draw[thick] (s1) -- (r) -- (s2);
  \node[font=\scriptsize,below=2mm of r,align=center] {uses IP addresses\\(L3)};
\end{tikzpicture}
\end{center}

\subsection{The firewall --- should every packet be allowed in?}
A router knows \emph{where} to send packets. But that raises a \textbf{security} question: should every packet be allowed to enter?
\begin{itemize}
  \item A \textbf{firewall} is a network security device.
  \item It \textbf{monitors} data coming into and going out of the network.
  \item Its main job: protect the network from unauthorised access or threats --- only authorised data gets through.
\end{itemize}

\textbf{Packet filtering.} A packet from outside passes the router, then the firewall inspects it:
\begin{center}
\begin{tikzpicture}
  \node[draw,fill=cnorange!25,minimum width=30mm,minimum height=9mm,align=center,font=\small] (ip) at (0,0) {IP header\\\scriptsize src / dst IP (L3)};
  \node[draw,fill=cngreen!25,minimum width=30mm,minimum height=9mm,align=center,font=\small,right=0mm of ip] (tcp) {TCP/UDP header\\\scriptsize src / dst port (L4)};
  \node[draw,fill=gray!15,minimum width=30mm,minimum height=9mm,align=center,font=\small,right=0mm of tcp] (d) {Data\\\scriptsize not examined};
  \draw[decorate,decoration={brace,amplitude=6pt}] (ip.north west) -- (tcp.north east) node[midway,above=7pt,font=\small] {simple packet filter checks these 4 values};
\end{tikzpicture}
\end{center}
So basic packet filtering decides on \textbf{four values}: source IP, destination IP, source port, destination port. IP is Layer 3, port is Layer 4 --- a packet-filtering firewall reads \textbf{two layers deep and stops}.

\begin{center}
\begin{tikzpicture}
  \node[draw,cloud,cloud puffs=10,fill=gray!10,minimum width=24mm,minimum height=14mm,font=\small] (i) at (0,0) {Internet};
  \node[rtr] (r) at (3,0) {R};
  \node[draw,fill=cnred!20,minimum width=12mm,minimum height=12mm,font=\small] (f) at (5.3,0) {Firewall};
  \node[draw,dashed,rounded corners,minimum width=26mm,minimum height=16mm,font=\small] (lan) at (8.4,0) {Internal network};
  \draw[arr] (i) -- (r); \draw[arr] (r) -- (f); \draw[arr,cngreen] (f) -- node[above,font=\scriptsize]{allowed} (lan);
  \draw[arr,cnred] (f) -- ++(0,-1.3) node[below,font=\scriptsize]{blocked};
\end{tikzpicture}
\end{center}

\subsection{Load balancers}
A \textbf{load balancer} distributes incoming traffic across \textbf{multiple servers}, so no single server is overwhelmed. In class, six client requests were spread across three servers instead of piling onto one. Netflix, Google, Amazon and Facebook all rely on load balancers.

\begin{center}
\begin{tikzpicture}
  \foreach \i in {1,...,6} \node[host,minimum width=9mm,minimum height=6mm,font=\scriptsize] (c\i) at (0,2.1-\i*0.6) {C\i};
  \node[draw,fill=cnpurple!20,minimum width=14mm,minimum height=16mm,align=center,font=\small] (lb) at (3.5,0) {Load\\balancer};
  \foreach \i in {1,...,6} \draw[->] (c\i.east) -- (lb.west);
  \foreach \i/\y in {1/1.2,2/0,3/-1.2} { \node[srv,font=\scriptsize] (s\i) at (7,\y) {Server \i}; \draw[arr] (lb) -- node[above,font=\scriptsize,sloped]{2 req.} (s\i); }
\end{tikzpicture}
\end{center}

\textbf{Link to Lecture 2:} it's the bottleneck idea again. When the slowest hop is a \emph{server} rather than a link, instead of making one server faster, you stop sending it everything.

\subsection{The devices by layer (memorise this!)}
\begin{center}
\begin{tabularx}{\textwidth}{>{\bfseries}l l X}
\toprule
Device & Layer & What it decides on \\
\midrule
Hub & L1 Physical & Nothing. Repeats signals to every port. \\
Switch & L2 Data Link & MAC addresses, within one LAN. \\
Router & L3 Network & IP addresses, between different networks. \\
Firewall & L3/L4 (and above) & IP addresses and port numbers; advanced types inspect more. \\
Load balancer & L4/L7 & Which of several servers should get this request. \\
Wireless router / access point & L1--L3 (by role) & Provide wireless access to the network. \\
\bottomrule
\end{tabularx}
\end{center}

\begin{center}
\begin{tikzpicture}[L/.style={draw,minimum width=30mm,minimum height=7mm,font=\small}]
  \node[L,fill=cnblue!25] (l7) at (0,0) {L7 Application};
  \node[L,fill=cngreen!25] (l4) at (0,-0.8) {L4 Transport};
  \node[L,fill=cnorange!25] (l3) at (0,-1.6) {L3 Network};
  \node[L,fill=cnpurple!20] (l2) at (0,-2.4) {L2 Data Link};
  \node[L,fill=gray!25] (l1) at (0,-3.2) {L1 Physical};
  \node[right=8mm of l1,font=\small] {Hub};
  \node[right=8mm of l2,font=\small] {Switch};
  \node[right=8mm of l3,font=\small] {Router};
  \node[right=30mm of l3,font=\small,yshift=4mm] {Firewall (L3--L4)};
  \node[right=8mm of l7,font=\small,yshift=-4mm] {Load balancer (L4--L7)};
  \draw[arr,cnred] (-2.2,-3.4) -- node[left,font=\scriptsize,align=right]{looks\\deeper} (-2.2,0.2);
\end{tikzpicture}
\end{center}

\section{Network topologies}
\textbf{Topology} = the arrangement of devices and their connections in a network, including how data flows. Types named: \textbf{star, mesh, hybrid, ring, bus}.

\begin{center}
\begin{tikzpicture}[n/.style={draw,circle,fill=cnblue!20,minimum size=4.5mm,inner sep=0pt}]
  % star
  \begin{scope}
    \node[draw,fill=cnorange!30,minimum size=6mm,font=\tiny] (c) at (0,0) {SW};
    \foreach \a [count=\i] in {90,162,234,306,18} {\node[n] (s\i) at (\a:1.2) {}; \draw (c) -- (s\i);}
    \node[font=\small\bfseries] at (0,-1.7) {Star};
  \end{scope}
  % mesh
  \begin{scope}[xshift=3.5cm]
    \foreach \a [count=\i] in {90,162,234,306,18} \node[n] (m\i) at (\a:1.1) {};
    \foreach \i in {1,...,5} \foreach \j in {1,...,5} { \ifnum\i<\j \draw (m\i) -- (m\j); \fi }
    \node[font=\small\bfseries] at (0,-1.7) {Mesh};
  \end{scope}
  % ring
  \begin{scope}[xshift=7cm]
    \foreach \a [count=\i] in {90,150,210,270,330,30} \node[n] (r\i) at (\a:1.1) {};
    \draw (r1)--(r2)--(r3)--(r4)--(r5)--(r6)--(r1);
    \node[font=\small\bfseries] at (0,-1.7) {Ring};
  \end{scope}
  % bus
  \begin{scope}[xshift=10.5cm]
    \draw[line width=2pt] (-1.5,0) -- (1.5,0);
    \foreach \x in {-1.2,-0.4,0.4,1.2} {\node[n] at (\x,0.8) {}; \draw (\x,0) -- (\x,0.6);}
    \foreach \x in {-0.8,0.8} {\node[n] at (\x,-0.8) {}; \draw (\x,0) -- (\x,-0.6);}
    \node[font=\small\bfseries] at (0,-1.7) {Bus};
  \end{scope}
\end{tikzpicture}
\end{center}

\subsection{Star topology}
\begin{itemize}
  \item All devices connect to a \textbf{central hub or switch}.
  \item Data flows sender $\to$ hub/switch $\to$ receiver.
  \item Centralised structure makes diagnostics and maintenance easy --- one obvious place to look.
  \item Popular for \textbf{LANs}; stays manageable as it grows.
  \item \textbf{Weakness:} the central device is a \textbf{single point of failure}.
\end{itemize}

\subsection{Mesh topology}
\begin{itemize}
  \item Every device is \textbf{directly connected to all others}.
  \item Many possible paths between any two devices $\Rightarrow$ high reliability: if one path fails, use another.
  \item Cost: the $n(n-1)/2$ problem --- jokingly called ``\textbf{spaghetti networks}''.
  \item Very useful when links are \textbf{wireless}: wireless mesh networks extend WiFi coverage, because adding a node adds paths, not cables.
\end{itemize}

\subsection{Hybrid topology}
Combines two or more topologies (e.g.\ star + ring + bus segments) into a more complex, flexible network. Real large networks are almost always hybrids because different parts of an organisation have different needs.

\subsection{Ring and bus (named in class)}
\begin{itemize}
  \item \textbf{Ring:} each device connects to two neighbours, forming a circle; data travels around the ring.
  \item \textbf{Bus:} all devices share one main cable (the ``bus'').
\end{itemize}

\subsection{Choosing a topology by environment}
\begin{center}
\begin{tabularx}{\textwidth}{l l X l}
\toprule
Environment & Topology & Reason & Example devices \\
\midrule
Small office & Star & Easy setup, centralised control & Laptops, printers, routers \\
Smart home & Mesh & Reliable, direct device-to-device communication & Smart bulbs, cameras \\
Data centre & Hybrid & Combines benefits of multiple topologies & Servers, storage units \\
\bottomrule
\end{tabularx}
\end{center}

\subsection{Fault tolerance and scalability}
\begin{center}
\begin{tabular}{lll}
\toprule
Topology & Fault tolerance & Scalability \\
\midrule
Star & Medium & Good \\
Mesh & High & Excellent \\
Hybrid & High & Hard \\
\bottomrule
\end{tabular}
\end{center}
Read this as \textbf{trade-offs}, not a ranking. ``Which topology is best?'' is not a good question; ask ``best for what, at what size, at what cost?''

\section{Hardware meets software: the cloud network}

\subsection{The evolutionary shift}
Traditional networking = physical boxes, special cables, manual configuration. In the cloud, the \textbf{logic is the same} (same layers, devices, topologies) but everything is \textbf{software-defined}.
\begin{center}
\begin{tabular}{lcl}
\toprule
Traditional & & Cloud era \\
\midrule
Physical box & $\to$ & Virtual instance \\
Copper cables & $\to$ & Virtualised overlays \\
Manual wiring & $\to$ & API calls or console clicks \\
\bottomrule
\end{tabular}
\end{center}
Nothing you learned becomes obsolete: a load balancer in AWS is still a load balancer; a subnet still partitions a Layer 3 network.

\subsection{AWS foundation: regions and edge locations}
\begin{itemize}
  \item \textbf{Regions:} isolated geographic areas (e.g.\ Mumbai), each with multiple data-centre clusters. 30+ globally, linked by an ultra-low-latency backbone.
  \item \textbf{Edge locations:} CDN endpoints that \textbf{cache content closer to users} to reduce latency.
\end{itemize}
Edge locations = Lecture 2 made concrete: caching closer shortens $d$ in $d/s$, so propagation delay falls. Same reason Netflix puts appliances inside ISPs.

\subsection{Redundancy: Availability Zones (AZs)}
\begin{itemize}
  \item An \textbf{AZ} = one or more discrete data centres with redundant power, networking and connectivity --- an isolated \textbf{fault domain}.
  \item Deploy across \textbf{two or more AZs} so your app survives a fire, flood or power failure at any one site.
\end{itemize}

\subsection{The Virtual Private Cloud (VPC)}
\begin{itemize}
  \item A \textbf{VPC} is your own private, software-defined network inside AWS. You control IP ranges, subnets and routing (\textbf{logical isolation}).
  \item \textbf{Subnets map to specific AZs}, partitioning your L3 network for scale and security.
  \item Addresses are written in \textbf{CIDR notation} (e.g.\ \texttt{10.0.0.0/16}): you get a block of addresses and decide how to carve it up.
\end{itemize}

\begin{center}
\begin{tikzpicture}
  \node[draw,thick,rounded corners,minimum width=120mm,minimum height=52mm,fill=cngrey] (reg) at (0,0) {};
  \node[anchor=north west,font=\small\bfseries] at (reg.north west) {Region (e.g.\ Mumbai)};
  \node[draw,thick,dashed,rounded corners,minimum width=110mm,minimum height=40mm,cnblue] (vpc) at (0,-0.3) {};
  \node[anchor=north west,font=\small\bfseries,cnblue] at (vpc.north west) {VPC  10.0.0.0/16};
  \node[draw,rounded corners,fill=cngreen!8,minimum width=48mm,minimum height=28mm] (az1) at (-2.7,-0.6) {};
  \node[anchor=north,font=\footnotesize\bfseries] at (az1.north) {Availability Zone A};
  \node[draw,fill=cnorange!15,font=\scriptsize,minimum width=38mm] at (-2.7,-0.6) {Subnet 10.0.1.0/24};
  \node[draw,rounded corners,fill=cngreen!8,minimum width=48mm,minimum height=28mm] (az2) at (2.7,-0.6) {};
  \node[anchor=north,font=\footnotesize\bfseries] at (az2.north) {Availability Zone B};
  \node[draw,fill=cnorange!15,font=\scriptsize,minimum width=38mm] at (2.7,-0.6) {Subnet 10.0.2.0/24};
\end{tikzpicture}
\end{center}

\begin{recap}
\begin{itemize}
  \item Full mesh needs $n(n-1)/2$ cables; a central device needs $n$.
  \item Hub = L1, broadcasts everything (unnecessary traffic). Switch = L2, MAC table, sends to correct port only. Router = L3, between networks.
  \item Firewall packet filter checks src/dst IP (L3) + src/dst port (L4), not data. Load balancer spreads requests over servers.
  \item Star (central, single point of failure), Mesh (reliable, costly ``spaghetti''), Hybrid (real world), Ring, Bus.
  \item Cloud = same ideas in software: Region, Edge location (CDN cache), AZ (fault domain), VPC (your private network), subnets per AZ, CIDR.
\end{itemize}
\end{recap}
